\chapter{机械能及其守恒}\label{ch:g11:mechanical-energy}

过了第一座小山，过山车便再没有发动机：每一个圈、
之后的每一次提速，都是从上坡时开设的“账户”中支取的。
本章把这笔账归结为两个量——动能与势能——以及一条规律：
只有重力做功时，二者之和保持不变。
摩擦力就是“手续费”，我们将学会从账目中
把它读出来。

\section{动能}

\begin{definition}[动能]\label{def:g11:mechanical-energy:kinetic}
质量为 $m$（\unit{kg}）、速率为 $v$（\unit{m/s}）的物体所具有的
\emph{动能}\index{动能}为
\[
E_k = \tfrac12\, m v^2 .
\]
单位检验：$\unit{kg.m^2/s^2} = \unit{N.m} = \unit{J}$——焦耳，
与其他能量一样（见\cref{ch:g10:energy-conservation}）。
\end{definition}

\begin{example}[数量级]\label{ex:g11:mechanical-energy:orders}
\begin{itemize}
  \item 一位行人，质量 \qty{70}{kg}，速率为 \qty{1.4}{m/s}：
        $E_k = \tfrac12 \times 70 \times 1.4^2 \approx \qty{69}{J}$;
  \item 一辆汽车，质量 \qty{1300}{kg}，速率为 \qty{130}{km/h} $= \qty{36.1}{m/s}$：
        $E_k \approx \qty{8.5e5}{J}$;
  \item 一列高速列车，$m = \qty{3.85e5}{kg}$，速率为 \qty{320}{km/h}
        $= \qty{88.9}{m/s}$: $E_k \approx \qty{1.5e9}{J}$;
  \item 一颗步枪子弹，质量 \qty{8.0}{g}，速率为 \qty{800}{m/s}：$E_k \approx \qty{2.6}{kJ}$。
\end{itemize}
七个数量级把散步与高速列车区分开来。而 $E_k$
随速率的\emph{平方}增长：在 \qty{130}{km/h} 时，汽车
携带的能量是 \qty{65}{km/h} 时的 $4$ 倍——
也正是刹车必须去除的能量的四倍。
\end{example}

\section{动能定理}

\begin{theorem}[动能定理]\label{thm:g11:mechanical-energy:ketheorem}
\index{动能定理}
对质量为 $m$ 的物体，沿直线从 $A$ 运动到 $B$，
其动能的改变量等于物体所受的
总功（见\cref{ch:g11:work-of-force}）：
\[
\Delta E_k = E_k(B) - E_k(A) = \sum W_{AB}(\vect F),
\]
式中求和遍历作用在物体上的所有力。该定理可推广
至沿任意路径的运动，这一点在本阶段不加证明地承认。
\end{theorem}
\admitted

\begin{example}[自由落体，逐帧检验]\label{ex:g11:mechanical-energy:freefall}
一个 \qty{0.200}{kg} 的小球被释放并拍摄下来；各帧数据如下：
\begin{center}\small
\begin{tabular}{lcccc}
$t$（\unit{s}）& 0.10 & 0.20 & 0.30 & 0.40\\
下落高度 $h$（\unit{m}）& 0.049 & 0.196 & 0.441 & 0.785\\
速率 $v$（\unit{m/s}）& 0.98 & 1.96 & 2.94 & 3.92\\
$E_k = \tfrac12 mv^2$（\unit{J}）& 0.096 & 0.384 & 0.864 & 1.54\\
$W = mgh$（\unit{J}）& 0.096 & 0.385 & 0.865 & 1.54
\end{tabular}
\end{center}
从 $E_k = 0$ 出发，每一帧的 $E_k$ 都与重力所做的功 $mgh$ 相等，
即 $\Delta E_k = W(\vect P)$，亦即 $v^2 = 2gh$——动能定理，
在读数精度范围内得到验证。
\end{example}

\begin{remark}[证明的出处]\label{rem:g11:mechanical-energy:proof}
一般结论——适用于任意路径，而不只是沿直线的
运动——可由牛顿第二定律及少量微积分推出，
两者均在本书后文中讲授（见\cref{ch:g12:newtons-laws,%
ch:g12:work-and-energy}）；在此之前，自由落体的检验便是我们的
依据；对曲线路径的版本，我们在本阶段不加证明地承认，并自由使用之。
\end{remark}

\section{重力势能}

\begin{definition}[重力势能]\label{def:g11:mechanical-energy:potential}
质量为 $m$ 的物体，其高度 $z$ 自选定的
\emph{参考平面}\index{参考平面}（$z = 0$ 处）向上度量，该物体所存储的
\emph{重力势能}\index{势能!重力}为
\[
E_p = m g z, \qquad g = \qty{9.81}{N/kg}.
\]
\end{definition}

\begin{remark}[只有差值有意义]\label{rem:g11:mechanical-energy:reference}
从 $A$ 运动到 $B$ 时，$E_p$ 的改变量为 $mg(z_B - z_A) =
-W_{AB}(\vect P)$：重力所做功的相反数
（见\cref{ch:g11:work-of-force}）。参考平面可自由选取——
地面、海平面、桌面皆可——这会使 $E_p$ 处处平移同一常量，
该常量在一切差值中都会抵消。把它放在数字最简单之处，
且在解题中途不要改动。
\end{remark}

\begin{example}[书架上的书]\label{ex:g11:mechanical-energy:shelf}
一本 \qty{1.2}{kg} 的书放在离地面 \qty{1.8}{m} 高处，而书架本身
高出街道 \qty{9.0}{m}。以地面为参考平面：
$E_p = 1.2 \times 9.81 \times 1.8 \approx \qty{21}{J}$；以街道为参考平面：
$E_p \approx \qty{127}{J}$。下落至地面时，
两种记账方式下释放的能量均为 \qty{21}{J}。
\end{example}

\section{机械能及其守恒}

\begin{definition}[机械能]\label{def:g11:mechanical-energy:em}
物体的\emph{机械能}\index{机械能}为
\[
E_m = E_k + E_p = \tfrac12 m v^2 + m g z .
\]
\end{definition}

\begin{theorem}[机械能守恒]\label{thm:g11:mechanical-energy:conservation}
\index{机械能守恒}
若物体只受其重力作用（无摩擦；其他
力，如无摩擦轨道的支持力，与
运动方向垂直），$E_m$ 守恒：对运动中的任意两点 $A$、$B$，
\[
\tfrac12 m v_A^2 + m g z_A = \tfrac12 m v_B^2 + m g z_B .
\]
\end{theorem}

\begin{proof}
由动能定理，$E_k(B) - E_k(A) = W_{AB}(\vect P) =
mg(z_A - z_B) = E_p(A) - E_p(B)$：$E_k$ 所增加的量等于 $E_p$ 所减少的量。
\end{proof}

\begin{example}[水上滑道]\label{ex:g11:mechanical-energy:slide}
一个小孩从无摩擦滑道顶端由静止开始下滑，滑道高
$h = \qty{3.2}{m}$。以底端为参考平面，$mgh = \tfrac12 mv^2$，于是
$v = \sqrt{2gh} \approx \qty{7.9}{m/s}$（\qty{29}{km/h}）——与
质量无关，\emph{也与滑道形状无关}：直线形、曲线形
或螺旋形，只有高度差起作用。
\end{example}

\begin{example}[单摆]\label{ex:g11:mechanical-energy:pendulum}
将单摆的摆球拉开，使其比最低点高出 $h = \qty{12}{cm}$，
然后释放。在两个端点处 $v = 0$；在最低点
$mgh$ 全部转化为动能，$v = \sqrt{2gh} \approx \qty{1.5}{m/s}$；
在另一侧，摆球恰好回升到 \qty{12}{cm} 高处。
\end{example}

\begin{omfigure}
\begin{tikzpicture}[scale=0.95, font=\small]
  \coordinate (O) at (0,0); \draw (-0.5,0) -- (0.5,0);
  \fill (O) circle (1.5pt);
  \draw[dashed, omExo] (230:2.2) arc (230:310:2.2);
  \draw (O) -- (230:2.2); \draw (O) -- (270:2.2); \draw (O) -- (310:2.2);
  \fill[omDef] (230:2.2) circle (2.5pt); \fill[omDef] (270:2.2)
    circle (2.5pt); \fill[omDef] (310:2.2) circle (2.5pt);
  \draw[-Stealth, omThm, thick] (270:2.2) -- ++(1.15,0)
    node[below] {$v_{\max}$};
  \node[above left] at (230:2.2) {$v = 0$};
  \node[above right] at (310:2.2) {$v = 0$};
  \draw[dashed, omExo] (230:2.2) -- ++(-0.9,0)
    (270:2.2) -- ++(-2.4,0);
  \draw[Stealth-Stealth] (-2.2,-2.2) -- (-2.2,-1.685)
    node[midway, left] {$h$};
\end{tikzpicture}

{\small 单摆在 $E_p$ 与 $E_k$ 之间相互转化：两端全为势能，
最低点全为动能，两侧升高的高度都是 $h$。}
\end{omfigure}

\begin{omfigure}
\begin{tikzpicture}[scale=0.75, font=\small]
  \draw[thick] plot [smooth, tension=0.7] coordinates
    {(0,3.2) (1.2,1.6) (2.4,0.5) (3.6,1.0) (4.6,2.0) (5.6,0.9)};
  \fill (0,3.2) circle (2pt) node[above right] {$A$};
  \fill (2.4,0.5) circle (2pt) node[below] {$B$}
    (4.6,2.0) circle (2pt) node[above] {$C$};
  \draw[omExo] (-0.2,0) -- (6.0,0);
  \fill[omProp!75] (7.0,0) rectangle (7.45,1.6);
  \fill[omProp!75] (7.9,0) rectangle (8.35,0.25)
    (8.8,0) rectangle (9.25,1.0);
  \fill[omDef!75] (7.9,0.25) rectangle (8.35,1.6)
    (8.8,1.0) rectangle (9.25,1.6);
  \draw[omExo] (6.7,0) -- (9.6,0);
  \draw[dashed, omThm] (6.7,1.6) -- (9.6,1.6) node[above left] {$E_m$};
  \node[below] at (7.225,0) {$A$}; \node[below] at (8.125,0) {$B$};
  \node[below] at (9.025,0) {$C$};
  \fill[omProp!75] (6.9,3.0) rectangle +(0.24,0.24);
  \fill[omDef!75] (6.9,2.5) rectangle +(0.24,0.24);
  \node[right] at (7.2,3.12) {$E_p$}; \node[right] at (7.2,2.62) {$E_k$};
\end{tikzpicture}

{\small 无摩擦过山车的能量柱状图：每个位置的
$E_p$ 柱与 $E_k$ 柱叠加后总量相同，均为 $E_m$（虚线）。}
\end{omfigure}

\begin{example}[跳台滑雪落地速率]\label{ex:g11:mechanical-energy:skijump}
一名跳台滑雪运动员以 $v_A = \qty{25}{m/s}$ 的速度离开起跳点，落点
比起跳点低 $\Delta z = \qty{40}{m}$。忽略空气阻力，
$v_B = \sqrt{v_A^2 + 2g\,\Delta z} = \sqrt{625 + 785} \approx
\qty{38}{m/s}$——无需知道飞行路径。
\end{example}

\begin{omfigure}
\begin{tikzpicture}[scale=0.75, font=\small]
  \draw[thick] (0,3.8) .. controls (1.2,2.7) and (1.9,2.4) .. (3.1,2.3);
  \draw[thick] (3.1,1.95) .. controls (4.2,1.55) and (5.5,1.0) .. (7.2,0);
  \draw[dashed, omThm, thick] (3.1,2.3) .. controls (4.3,2.25) and
    (5.2,1.7) .. (5.9,0.75);
  \fill (3.1,2.3) circle (2pt) node[above] {$A$}
    (5.9,0.75) circle (2pt) node[above right] {$B$};
  \draw[dashed, omExo] (3.1,2.3) -- (8.2,2.3) (5.9,0.75) -- (8.2,0.75);
  \draw[Stealth-Stealth] (8.0,0.75) -- (8.0,2.3)
    node[midway, right] {$\Delta z$};
\end{tikzpicture}

{\small 跳台滑雪：无论从 $A$ 到 $B$ 飞行的是怎样的曲线，只有
高度差 $\Delta z$ 进入能量关系式。}
\end{omfigure}

\section{摩擦耗散}

\begin{definition}[耗散]\label{def:g11:mechanical-energy:dissipation}
离开机械能的能量称为
\emph{耗散}\index{耗散}：它以热能的形式
重新出现在相互摩擦的表面与空气中。总能量是守恒的
（见\cref{ch:g10:energy-conservation}）；而机械能本身则不然。
\end{definition}

\begin{proposition}[亏损量衡量摩擦力]\label{prop:g11:mechanical-energy:deficit}
若摩擦力 $\vect f$ 也沿运动方向做功，
\[
\Delta E_m = E_m(B) - E_m(A) = W_{AB}(\vect f) < 0:
\]
机械能的损失等于摩擦力所做的功。
\end{proposition}

\begin{proof}
动能定理此时写为 $\Delta E_k = W_{AB}(\vect P) +
W_{AB}(\vect f) = -\Delta E_p + W_{AB}(\vect f)$；将 $\Delta E_p$
移至另一侧。摩擦力阻碍运动，所以其做功为负。
\end{proof}

\begin{example}[雪橇能量核算]\label{ex:g11:mechanical-energy:toboggan}
一名儿童与雪橇（$m = \qty{40}{kg}$）从静止开始，
沿 \qty{12}{m} 长的积雪路段下降 $h = \qty{5.0}{m}$，末速率为
\qty{7.0}{m/s}，而无摩擦时应为
$\sqrt{2gh} \approx \qty{9.9}{m/s}$。取坡底为参考平面：
$\Delta E_m = \tfrac12 \times 40 \times 7.0^2 - 40 \times 9.81
\times 5.0 = 980 - 1962 = \qty{-982}{J}$。因此 $W(\vect f) =
\qty{-982}{J}$，路程 \qty{12}{m} 上的平均摩擦力
$f = 982/12 \approx \qty{82}{N}$，另有 \qty{982}{J} 的
热量留在滑板与雪地中。
\end{example}

\begin{omfigure}
\begin{tikzpicture}
\begin{axis}[omaxis, width=8.6cm, height=4.6cm,
    xlabel={$t$ (\unit{s})}, ylabel={$E_m$ (\unit{J})},
    xmin=0, xmax=10, ymin=0, ymax=1250, legend pos=south west]
  \addplot[omThm, thick, domain=0:10, samples=2] {1000};
  \addplot[omDef, thick, domain=0:10, samples=60] {1000*exp(-0.18*x)};
  \legend{无摩擦，有摩擦}
\end{axis}
\end{tikzpicture}

{\small 沿运动的 $E_m$：无摩擦时保持水平，有摩擦时随之
逐渐减小——每一时刻的下降量即此前产生的热量。}
\end{omfigure}

\begin{method}[能量核算]\label{met:g11:mechanical-energy:bookkeeping}
\begin{enumerate}
  \item 选定参考平面并保持不变；记 $A$（已知状态）、$B$
        （所求状态）。
  \item 写出 $E_m = \tfrac12 mv^2 + mgz$，并分别用于 $A$ 与 $B$ 处。
  \item 无摩擦时 $E_m(A) = E_m(B)$；求解——通常为
        $v_B = \sqrt{v_A^2 + 2g(z_A - z_B)}$。
  \item 有摩擦时 $E_m(B) - E_m(A) = -f\ell$（$\ell$ 为路径长）；
        解出未知量。
  \item 校验：$E_k$ 与 $E_p$ 的条形叠加为 $E_m$，除非有摩擦，否则保持水平。
\end{enumerate}
\end{method}

\section{练习}

\begin{exercise}[$\star$]\label{exo:g11:mechanical-energy:1}
求下列各物体的动能：(a) 质量为 \qty{90}{kg} 的踏板车与骑车人，以
\qty{25}{km/h} 行驶；(b) 质量为 \qty{1200}{kg} 的汽车以 \qty{130}{km/h} 行驶；(c) 速率加倍时，
$E_k$ 变为原来的多少倍？速率变为三倍时又如何？
\end{exercise}

\begin{exercise}[$\star$]\label{exo:g11:mechanical-energy:2}
一名质量为 \qty{75}{kg} 的攀登者上升了 \qty{850}{m} 的高度，求其
势能的增量。它是否取决于参考平面？是否取决于路径？
\end{exercise}

\begin{exercise}[$\star$]\label{exo:g11:mechanical-energy:3}
一只花盆从阳台坠落 \qty{20}{m}。忽略空气阻力，
求其落地速率（分别以 \unit{m/s} 与 \unit{km/h} 表示）。质量哪里去了？
\end{exercise}

\begin{exercise}[$\star$]\label{exo:g11:mechanical-energy:4}
一辆质量为 \qty{1200}{kg} 的汽车以 \qty{25}{m/s} 的速率刹车，经 \qty{50}{m} 路程后静止。
用动能定理求制动力所做的功，再求
平均制动力。
\end{exercise}

\begin{exercise}[$\star$]\label{exo:g11:mechanical-energy:5}
一单摆摆球从高于最低点 \qty{8.0}{cm} 处释放。求
其在最低点的速率，以及摆至另一侧所达到的高度。
\end{exercise}

\begin{exercise}[$\star\star$]\label{exo:g11:mechanical-energy:6}
利用\cref{ex:g11:mechanical-energy:orders}，火车的动能是子弹动能的
多少倍？火车自身的 $E_k$ 能将其举升至多高
（$h = v^2/2g$）？
\end{exercise}

\begin{exercise}[$\star\star$]\label{exo:g11:mechanical-energy:7}
两条无摩擦的滑道从同一 \qty{3.0}{m} 高的平台滑下，一条
笔直而陡峭，一条长而蜿蜒。比较二者到达底端的速率
（需计算）与到达时间（无需计算）。
\end{exercise}

\begin{exercise}[$\star\star$]\label{exo:g11:mechanical-energy:8}
一小球以 \qty{12}{m/s} 的速率竖直上抛。求其能达到的最大高度，
再求经过 \qty{4.0}{m} 高度处的速率——上升与下落两种情况下的
速率。
\end{exercise}

\begin{exercise}[$\star\star$]\label{exo:g11:mechanical-energy:9}
一名滑雪运动员以 \qty{23}{m/s} 的速率离开起跳台，落点比起跳台
低 \qty{35}{m}。不计空气阻力，预测落地速率；实际值是
偏大还是偏小？并说明理由。
\end{exercise}

\begin{exercise}[$\star\star$]\label{exo:g11:mechanical-energy:10}
一质量为 \qty{45}{kg} 的雪橇由静止开始，沿 \qty{15}{m} 长的斜坡
下降 \qty{6.0}{m}，到达底端时速率为 \qty{8.0}{m/s}。求损失的
机械能，再求平均摩擦力。能量去了哪里？
\end{exercise}

\begin{exercise}[$\star\star$]\label{exo:g11:mechanical-energy:11}
伽利略的钉子：摆锤从最低点上方 \qty{15}{cm} 处释放，
摆线竖直时，遇到一颗挡住摆线上半段的钉子。求
摆锤到达最低点时的速率，以及越过钉子后上升的高度；钉子改变了什么，
又不改变什么？
\end{exercise}

\begin{exercise}[$\star\star\star$]\label{exo:g11:mechanical-energy:12}
一名滑板运动员从 \qty{2.5}{m} 高处滑入半管，每次横越
（来回一边到另一边）耗散 $10\%$ 的机械能。求
第一次下滑到底部时的速率，再通过计算连续上升的高度，
求首次使上升高度低于 \qty{1.0}{m} 的横越。
\end{exercise}

\begin{exercise}[$\star\star\star$]\label{exo:g11:mechanical-energy:13}
从 \qty{45}{m} 高的悬崖顶端，以 \qty{15}{m/s} 的速率抛出一块石头——
可以竖直向上抛、水平抛，或任意角度抛出。证明各种情形下落地速率相同，
并求出该速率。究竟什么\emph{确实}与角度有关？
\end{exercise}

\begin{exercise}[$\star\star\star$]\label{exo:g11:mechanical-energy:14}
一名质量为 \qty{80}{kg} 的滑雪者由静止出发，
沿 \qty{800}{m} 长的滑道下滑，高度下降 \qty{120}{m}，受恒定的 \qty{65}{N}
摩擦力作用。求到达底端的速率，以及
被耗散的 $E_m$ 所占的比例。
\end{exercise}

\begin{exercise}[$\star\star\star$]\label{exo:g11:mechanical-energy:15}
一座抽水蓄能电站将 $V = \qty{2.0e6}{m^3}$ 的水
（\qty{1000}{kg} 每 \unit{m^3}）提升 \qty{300}{m} 高度。计算
蓄存的能量，分别以焦耳和 \unit{kW.h} 表示；水流经水轮机放能时效率为 $85\%$
（见\cref{ch:g10:energy-conservation}），这些电能按每户每天
用电 \qty{10}{kW.h} 计，可供多少户家庭用电一天？
\end{exercise}

\section{问题：设计环形轨道}

\begin{problem}[{周末问题——过山车的能量审计：
圆环规定的发射高度、它许诺的各处速率、
两个传感器坦白的摩擦、结清账目的刹车——
整段行程系于同一个数}]
\label{pb:g11:mechanical-energy:1}
一列质量为 $m = \qty{2.0e3}{kg}$ 的列车，由静止从发射坡上
高 $H$ 处释放，无动力滑行，穿过竖直的圆环轨道，
圆环半径 $R = \qty{8.0}{m}$，顶点位于高 $2R$ 处。
圆周运动的分析（下学年学习）提供了我们借用的一个数据：在
圆环顶点处，列车压向轨道的条件是
$v_{\text{顶}}^2 \geq gR$。I 部分和 II 部分忽略摩擦。

\medskip
\noindent\textbf{I 部分——圆环所要求的坡高。}
\begin{enumerate}
  \item 求圆环顶点处的最小安全速率。
  \item 以地面为参考平面，求列车
        在圆环顶点、速率取最小值时的机械能。
  \item 推出 $H_{\min}$，并证明 $H_{\min} = 5R/2$——结果中既无质量，也无 $g$。
  \item 车队取 $H = \qty{24}{m}$，求
        圆环顶点的实际速率，以及比值 $v_{\text{顶}}^2/(gR)$。
  \item $m$ 去了哪里？为什么这一设计
        对满载与空载的列车都适用？
\end{enumerate}

\medskip
\noindent\textbf{II 部分——高度许诺了什么}（$H = \qty{24}{m}$）。
\begin{enumerate}[resume]
  \item 圆环底部（$z = 0$）处的速率，以 \unit{m/s} 和
        \unit{km/h} 为单位各是多少？
  \item 圆环侧面（$z = R$）处的速率？
  \item 计算发射点及圆环底部、侧面、顶部处的 $E_p$ 与 $E_k$，
        并验证每一组之和都等于 $mgH$。
  \item 圆环之后接一个高 \qty{12}{m} 的驼峰：峰顶速率是多少？
  \item 给出 $v(z)$ 的一般表达式；过山车在何处
        速率最大，又在何处最小？
\end{enumerate}

\medskip
\noindent\textbf{III 部分——摩擦力审计。}
两个传感器位于同一高度 $z = \qty{2.0}{m}$，沿轨道
相距 \qty{60}{m}，依次读得 $v_1 = \qty{20.4}{m/s}$ 与 $v_2 = \qty{19.2}{m/s}$。
\begin{enumerate}[resume]
  \item 计算两个传感器之间机械能的变化量。
        这种同高度安置为什么巧妙？
  \item 推出摩擦力所做的功及平均摩擦力。
  \item 缺失的能量去了哪里？说出它的形式及去向。
  \item 每经过 \qty{60}{m}，列车损失的 $E_m$ 占多大比例？为什么
        为何仅取 $H_{\min}$ 在真实的、有摩擦的世界中不安全？
  \item 发射点与圆环顶点之间约有 \qty{90}{m} 的轨道：估算
        那里的 $E_m$；第 4 题中的余量是否仍能保住圆环？
\end{enumerate}

\medskip
\noindent\textbf{IV 部分——结清账目。}
最后一个传感器位于水平直线制动段的入口
（$z = 0$）处，读数为 \qty{18.0}{m/s}。
\begin{enumerate}[resume]
  \item 计算列车在该处的动能。
  \item 制动须在 $d = \qty{45}{m}$ 内使列车停下，求
        所需的（恒定）制动力。
  \item 舒适性检验：计算减速度 $a = F/m$，并与 $g$ 比较。
  \item 整段行程中，轨道摩擦耗散了多少能量？
        验证摩擦 $+$ 制动 $= mgH$，精确到焦耳。
  \item 结语（用一句话）：哪一个数决定了圆环的
        安全性、各处的速率、摩擦力造成的损耗和制动力——
        以及忽略摩擦所付出的代价是什么？
\end{enumerate}
\end{problem}
